%%%% PROCESAR con PdfLaTeX !!!!!


\documentclass[12pt]{book}
\usepackage{geometry}\geometry{top=2cm,bottom=2cm,left=3cm,right=3cm}
\usepackage{amssymb}
\usepackage{amsmath}
\usepackage{graphicx}
\usepackage{txfonts}

\usepackage[hidelinks]{hyperref}
\usepackage[spanish]{babel}
\setcounter{tocdepth}{3}
\usepackage[usenames]{color}

\usepackage{framed}
\usepackage{color}
\usepackage{wrapfig}\definecolor{shadecolor}{RGB}{224,238,238}

\usepackage[hidelinks]{hyperref}
\hypersetup{
    colorlinks=true,
    linkcolor=blue,
    filecolor=magenta,      
    urlcolor=cyan,
    pdftitle={Sharelatex Example},
    bookmarks=true,
%    pdfpagemode=FullScreen,
}

\begin{document}
\thispagestyle{empty}

\begin {center}

\includegraphics[scale=.4]{Logo-fiuba_big.png}

\medskip
UNIVERSIDAD DE BUENOS AIRES

Facultad de Ciencias Ex\'actas y Naturales

Departamento de Matem\'aticas


\vspace{3cm}


\textbf{\large ÁLGEBRA I}

\vspace{2cm}


Este es un modesto aporte para los alumnos de la f\'acultad de Ciencias Exactas y Naturales de la UBA de las carreras de licenciatura en Matemática y Computaci\'on.
De ninguna man\'era pretende ser una gu\'ia de estudio, ni remplaza las clases presenciales, el material oficial de la catedra esta disponible en el web site de la m\'ateria.
\\
http://cms.dm.uba.ar/

\end {center}


\vspace{2.5cm}

\noindent Autor:\,	Isaac Edgar Camacho Ocampo
 
\noindent Carrera:\,	Licenciatura en Ciencias de la Computaci\'on

\vspace{1cm}

\vspace{1cm}

\noindent Buenos Aires, 2020
\\ 
\\
\textit{Si se encuentra alg\'un error u omisi\'on en este res\'umen por favor colaborar en \\
\url{https://github.com/IsaacEdgarCamacho/Apuntes/tree/master/Algebra/} \quad \\ o escribirme a \\ \url{isaac.edgar.camacho@gmail.com}
}

\newpage


\tableofcontents

\tableofcontents
\chapter{Introducción}
\section{Conocimientos previos}
\section{Estado del arte}


\chapter{Conjuntos, Relaciones y Funciones.}
\section{Conjuntos.}
\subsection{Conjuntos y subconjuntos, pertenencia e inclusión.}
\subsubsection{Definición 1.1.1. (informal de conjunto y elementos.)}
Un conjunto es una colección de objetos, llamados elementos, que tiene la
propiedad que dado un objeto cualquiera, se puede decidir si ese objeto es
un elemento del conjunto o no.
\\
\\
\textit{Ejemplos:}
\begin{itemize}
\item $A = \{1, 2, 3\} ,\quad B = \{\triangle, \square, \lozenge \} , \quad  C = \{1, \{1\}, \{2, 3\}\} .$
\item $N = \{1, 2, 3, 4, \dots \} \quad $ el conjunto de los números naturales.
\end{itemize}


\subsection{Operaciones entre conjuntos}

\begin{shaded}
\textbf{Teorema }:si un conjunto A esta incluido en otro conjunto B, entonces el complemento de B esta incluido en el complemento de A
\end{shaded}
\begin{center}$ A\subseteq B \Rightarrow  B^{c} \subseteq A^{c}$\end{center}


\textit{Demostracion:}

\begin{equation*}
\begin{split}
Hitpotesis:   \quad  A \subseteq B \longleftrightarrow \quad \forall x \epsilon A \Rightarrow x \epsilon B  \\
&\equiv  \quad  def( a) \wedge def(i) \wedge \ ( 0 \leq i < \vert a \vert ) \wedge (0 \leq i+1 < \vert a \vert ) \quad      \wedge_{L} \ True  \\ \\
wp(S, \ Post) \ &\equiv  \quad  ( 0 \leq i < \vert a \vert ) \wedge (0 \leq i+1 < \vert a \vert )
\end{split}
\end{equation*}

\begin{equation*}
\begin{split}
wp(S, \ Post) \  &\equiv \quad wp ( \ result  := a[i] + a[i+1] \ , \ Post^{result}_{a[i] + a[i+1]} \ )  \\ \\
&\equiv  \quad  def( a[i] + a[i+1]) \quad \wedge_{L} \quad \ \underbrace{a[i] + a[i+1] = a[i] + a[i+1]}_{\text{es siempre True}}  \\
&\equiv  \quad  def( a) \wedge def(i) \wedge \ ( 0 \leq i < \vert a \vert ) \wedge (0 \leq i+1 < \vert a \vert ) \quad      \wedge_{L} \ True  \\ \\
wp(S, \ Post) \ &\equiv  \quad  ( 0 \leq i < \vert a \vert ) \wedge (0 \leq i+1 < \vert a \vert )
\end{split}
\end{equation*}

El Teorema de Pitágoras lleva este nombre porque su descubrimiento recae sobre la escuela pitagórica. (Fuente: Wikipedia)


\begin{shaded}
\textbf{Ejercicio: }
Probar usando el principio de inducción que 7 divide a $ 15^{n} + 7$ para todo $n \in \mathbb N $  
\end{shaded}

\begin{wrapfigure}[0]{r}{0.5\textwidth}
\vspace{-15pt}
\begin{shaded}
\textbf{Ejercicio: } Probar usando el principio de inducción que 7 divide a $ 15^{n} + 7$ para todo $n \in \mathbb N $  
\end{shaded}
\end{wrapfigure}

\begin{equation*}
\begin{split}
7\mid (15^{n} + 7) \  &\equiv \quad wp ( \ result  := a[i] + a[i+1] \ , \ Post^{result}_{a[i] + a[i+1]} \ )  \\ \\
&\equiv  \quad  def( a[i] + a[i+1]) \quad \wedge_{L} \quad \ \underbrace{a[i] + a[i+1] = a[i] + a[i+1]}_{\text{es siempre True}}  \\
&\equiv  \quad  def( a) \wedge def(i) \wedge \ ( 0 \leq i < \vert a \vert ) \wedge (0 \leq i+1 < \vert a \vert ) \quad      \wedge_{L} \ True  \\ \\
wp(S, \ Post) \ &\equiv  \quad  ( 0 \leq i < \vert a \vert ) \wedge (0 \leq i+1 < \vert a \vert )
\end{split}
\end{equation*}




\section{Ejercicios}


\begin{wrapfigure}[0]{r}{1\textwidth}
\vspace{-15pt}
\begin{shaded}
\textbf{Ejercicio: } Probar usando el principio de inducción que 7 divide a $ 15^{n} + 7$ para todo $n \in \mathbb N $  
\end{shaded}
\end{wrapfigure}


\chapter{Resultados}
\section{Simulación de resultados}
\subsection{Suposiciones}
\subsection{Modelos}
\section{Resultados preliminares}
\section{Resultados postprocesados}
\subsection{Valores atípicos}
\subsection{Correlaciones}
\chapter{Conclusiones}


\begin{shaded}
\textbf{Ejercicio:} Probar usando el principio de inducción que 7 divide a $ 15^{n} + 6$ para todo $n \in \mathbb N_{0}$  
\end{shaded}

Cuando trabajamos con inducci\'on debemos definir cual es la afirmacion que vale para cada n   y no se pone para todo n, porque eso es lo que queremos probar.\\
\[
 P(n): 7 \mid 15^{n} + 6
\]
\newline
\textbf{Caso base $P(0): n = 0  \ $}\quad Queremos ver si el primer numero del conjunto cumpler P(n)

\[
(n=0 ) \ \Rightarrow \ 7 \ divide \ a \ (15^{0} + 6 )
\]

\[  15^{0} + 6 = 7  , \  7 = 1 \cdot 7  \ \Rightarrow 7\mid 7 \]
\\
Por lo tanto la afirmaci\'on para $P(0)$  es verdadera ya que 7 divide a 7.
\newline
\newline
\textbf{Paso inductivo} Suponemos que $P(n)$ es verdadera y probamos la veracidad de $P(n+1)$
\[
\text {Hip\'otesis  Induct\'iva } \ (HI) \equiv \ (15^{n} + 6 = k \cdot 7) ,\ con \ k \in \mathbb{Z} 
\]
\newline
Lo que queremos hacer es que suponiendo que $P(n)$ es verdad probamos la verdad de $P(n+1)$

\begin{equation*}
\begin{split}
15^{n+1} + 6 \quad   &\equiv \quad  15^{n+1} + 6   \\
&\equiv  \quad   15^{n} \cdot 15 + 6  \\
&\equiv  \quad   15^{n} \cdot (1 + 14) + 6 \\
&\equiv  \quad   15^{n} \cdot 1 + 15^{n} \cdot 14 + 6 \\ 
&\equiv  \quad   (15^{n} \cdot 1 + 6) +15^{n} \cdot 14  \\
HI \quad   &\equiv  \quad   k \cdot 7 +15^{n} \cdot 14  \\
\quad   &\equiv  \quad   k \cdot 7 +15^{n} \cdot (2 \cdot 7)  \\
15^{n+1} + 6 \quad   &\equiv  \quad  7 \cdot \underbrace{ (k +15^{n} \cdot 2)}_{ esto \ \in \mathbb{Z}}  \\
\end{split}
\end{equation*}
Podemos ver que $ 15^{n+1} + 6  $ se expresa como multiplo de 7, por lo tanto queda demostrado nuestro ejercicio

\begin{shaded}
\textbf{Ejercicio:} Probar usando el principio de inducción que 10 divide a $ 16^{n} - 6^{n}$ \\ para todo $n \in \mathbb N$  \\
\[	
	10  \mid 16^{n} - 6^{n} \quad \forall n \in \mathbb{Z}
\]
\end{shaded}

\begin{equation*}
\begin{split}
\textbf{Caso Base P(1) } & \quad n=1 \\
16^{1} - 6^{1} \quad   &\equiv \quad  16 - 6 = 10 \\
&\equiv \quad  10 = 1 \cdot 10 \text{\quad (con esto vemos que cumple)}\\ \\
\textbf{Paso inductivo P(n)} & \quad P(n) \text { es verdad y pruebo para } P(n+1)	\\
\text{Hip\'otesis inductiva } &\Rightarrow \quad (16^{n} - 6^{n} = k \cdot 10) \text{ con } k \in \mathbb{Z}\\ \\
16^{n+1} - 6^{n+1} \quad &\equiv  \quad   16^{n} \cdot 16 - 6^{n} \cdot 6  \\
&\equiv  \quad   16^{n} \cdot (6+10) - 6^{n} \cdot 6  \\
&\equiv  \quad   (16^{n} \cdot 6 + 16^{n} \cdot10) - 6^{n} \cdot 6  \\
&\equiv  \quad    16^{n} \cdot 6 - 6^{n} \cdot 6  +  16^{n} \cdot 10 \\
&\equiv  \quad   6 \cdot \underbrace{ (16^{n} - 6^{n}) }_{HI}  +  16^{n} \cdot 10 \\
HI \quad   &\equiv  \quad  6 \cdot k \cdot 10 + 16^{n} \cdot 10 \\
16^{n+1} - 6^{n+1} \quad  &\equiv  \quad  10 \cdot (\underbrace{  k \cdot 6 + 16^{n})}_{ esto \ \in \mathbb{Z}}  
\end{split}
\end{equation*}
Podemos ver que $ 16^{n+1} - 6^{n+1} $ se expresa como multiplo de 10, por lo tanto queda demostrado nuestro ejercicio

\end{document}
